Supernova SN 1987A was at its brightest with an apparent magnitude of $+3$ on
about 15th May 1987 and then faded, finally becoming invisible to the naked
eye by 4th February 1988. It is assumed that brightness $B$ varied with time
$t$ as an exponential decline, $B = B_0 e^{-t/\tau}$, where $B_0$ and $\tau$
are constant. The maximum apparent magnitude which can be seen by the naked
eye is $+6$.

\begin{description}
\item[a)] Determine the value of $\tau$ in days. \hfill[5]
\item[b)] Find the last day that observers could have seen the supernova if
  they had a 6-inch (15.24-cm) telescope with transmission efficiency
  $T = 70\%$. Assume that the average diameter of the human pupil is
  0.6\,cm. \hfill[10]
\end{description}
